\documentclass[10pt]{article}

\usepackage[utf8]{inputenc}

\usepackage[T1]{fontenc}

\usepackage{amsmath}

\usepackage{amsfonts}

\usepackage{amssymb}

\usepackage[version=4]{mhchem}

\usepackage{stmaryrd}

\usepackage{bbold}

\usepackage{graphicx}

\usepackage[export]{adjustbox}

\graphicspath{ {./images/} }



\begin{document}

в нек-ром «основном» пространстве $X$; геометрич. процессы определяются как точечные процессы на многообразиях, представляющих пространство соответствующих элементов. Так, процессы прямых на плоскости определяются как точечные процессы на листе Мёбиуса (последний представляет пространство прямых на $\mathbb{R}^{2}$ ). Рассматриваются также процессы $d$-мерных плоскостей в $\mathbb{R}^{n}$, процессы выпуклых фигур в $\mathbb{R} n$, случайные мозаики в $\mathbb{R}^{n}$ (последние можно рассматривать как процессы выпуклых многогранников, в реализациях к-рых с вероятностью 1 внутренности многогранников не пересекаются, а объединение их замыканий дают все $\mathbb{R}^{n}$ ) и т. Д.



Более общим понятием являются процессы многообразий; здесь С. г. смыкается с теорией случайных множеств (см. [1]).



Другой особенностью, выделяющей С. г. из теории случайных множеств, является пнтерес С. Г. к геометрич. процессам с распределениями, инвариантными относительно грушп, действующих в основном пространстве $X$. Характерным результатом в этом направлении является следующий (см. [2]). Рассматривается класс процессов прямых на $\mathbb{R}^{2}$, обладающих конечной интенсивностью и инвариантных относительно группы евклидовых движений плоскости. Процесс наз. н е о с об е н в ы м, если его распределение Палма абсолютно непрерывно относительно (безусловного) распределения процесса. Все неособенные процессы прямых суть дважды стохастические пуассоновские (т. е. пуассоновские, управляемые случайной мерой). Для точечных полей в $\mathbb{R}^{n}$ это не имеет места.



Ряд столь же неожиданных свойств других геометрич. процессов, инвариантных относительно групп, было обнаружено с помощью аппарата комбинаторной интегральной геометрии (см. [3]). Способом усреднения комбинаторных разложений по простравству реализаций процесса был получен, напр., следующий результат (см. [3]). Рассматривается случайное множество $U$ на $\mathbb{R}^{2}$, представляющее собой объединение областей из нек-рого процесса выпуклых областей, инвариантного относительно евклидовой группы. Пусть $U-$ черное множество, его дополнение - белое. Альтернирующий процесс черных и белых интервалов, индуцируөмый множеством $U$ на прямой $O x$, наз. черно-рекуррентным, если:



a) белые интервалы $a_{i}$ составляют независимый процесс восстановления;



б) тройки $\left(b_{i}, a_{i}, \beta_{i}\right)$, где $b_{i}$ - длина $i$-го черного интервала, $a_{i}, \beta_{i}$ - углы пересечения $O x$ с $\partial U$ на концах $b_{i}$, для разных $i$ независимы.



При общих предположениях эргодичности и существования нек-рых моментов, а также при отсутствии прямолинейных участков на $\partial U$, из черно-рекуррентности следует, что длина белого интервала распределена әкспоненциально.



Важное значение в С. г. имеет развитое здесь понятие «типичного» элемента данного геометрич. процесса. Рассматриваются задачи об ошисании распределений, удовлетворяющих различным условиям «типичных» $k$-подмножеств элементов в геометрич. процессах (пример такой задачи: найти распределение евклидово инвариантных характеристик «типичного» треугольника с вершинами в реализациях точечного процесса, причем требуется, чтобы внутренность треугольника содержала бы $l$ точек из реализации). Решение таких задач получено для пуассоновских процессов. Подобные задачи возникают, напр., в астрофизике.



Задачи т. н. стереологии также относятся к С. г., если они ставятся для процессов геометрич. фигур. (В стереологии требуется описать многомерный образ по его сечениям прямыми или плоскостями меньшего тисла измерений.) Здесь имеются результаты по стереологии первой и второй моментной меры. Выше были затронуты лишь наиболее характерные задачи, т. к. границы С. г. вряд ли можно точно определить. К С. г. примыкают следующие области: геометрич. статистика [4], теория (случайных) множеств дробной размерности [5], математич. морфология и анализ изображений $[6]$.



Лит.: [1] М а т е р о н Ж., Случайные множества и инте-



\includegraphics[max width=\textwidth, center]{2023_07_09_5f04b0ebf4a2525306b6g-1}

d ing E. F., K e n d a l l D. G., Stochastic geometry, N. Y., geometry, N. Y., 1982; [4] R i p l e y B. D., Spatial statistics, N. Y.-[a. 0.], 1981; [5] M a n d e l b r o t B. B., Fractals: Form, chance and dimension, S. F., 1977; [6] S e r r a J., Image analysis and mathematical morphology, N. Y., 1982 .



СТОХАСТИЧЕСКАЯ ЗАВИСИМОСТЬ (в е р о я Тност в а я, с т а т и с т и ч е к а я) - зависимость между случайными величинами, к-рая выражается в изменении условных распределений любой из величин при изменении значений других величин. Виды С. 3. многообразны: если случайные величины не являются взаимно независимыми, то им в той или иной степени свойственна С. з. Одним из наиболее общих типов С. з. является корреляционная зависимость (см. Корреляция, Регрессия). Из конкретных видов С. З. наиболее изучена марковская зависимость (см. Маркова цепь, Марковский процесс, Марковское свойство).



См., кроме того, Случайный процесс, Стационарный случайный процесс. $\quad$ A. В. Прохоров.



СТОХАСТИЧЕСКАЯ ИГРА - динамическая иәра, у к-рой переходная функция распределения не зависит от предыстории игры, т. е.



$F\left(x_{k} \mid x_{1}, s^{\left(x_{1}\right)}, \ldots, x_{k-1}, s^{\left(x_{k-1}\right)}\right)=F\left(x_{k} \mid x_{k-1}, s^{\left(x_{k-1}\right)}\right)$.



С. и. были впервые определены Л. Шепли [1], к-рый рассматривал антагонистич. С. и. с интегральным выигрышем (и г р ы Ш е п л и). В играх Шепли как множество $X$ состояний игры, так и множества әлементарных стратегий игроков конечны и, кроме того, на любом шаге при любом выборе игроками альтернатив имеется ненулевая вероятность окончания партии. Вследствие последнего условия, партия с вероятностью 1 заканчивается за конечное число шагов, и математич. ожидание выигрыша каждого из игроков конечно. Любая такая игра обладает значением и оба игрока имеют с т а ц и он а р пы е оптимальны е с т р а тетии, т. е. стратегии, в к-рых выбор игроком әлементарной стратегии в каждом состоянии игры зависит лишь от текущего состояния. Им же была указана процедура, дающая возможность найти как значение игры, так и оптимальные стратегии.



Рассматривались также С. и., отличающиеся от игр IIепли возможностью бесконечных партий, с п р ед ельны м с р е дним выи г рышем, т. е. антагонистич. С. и. с



\[

h_{1}(P)=-h_{2}(P)=\lim _{n \rightarrow \infty} \sup \sum_{k=1}^{n} \frac{1}{n} h_{i}\left(x_{k}, s^{\left(x_{k}\right)}\right) .

\]



Было показано существование значения такой игры и стационарных оптимальных стратегий в предположении эргодичности марковской цепи, возникающей при подстановке в переходные функции $F\left(x_{k} \mid x_{k-1}, s\left({ }^{x_{k}-1}\right)\right)$ любых стационарных стратегий. Эти результаты обобщались как в направлении снятия ограничений на число состояний и әлементарных стратегий, так и на случай иных форм выигрышей.



Jum.: [1] S $\mathrm{h}$ a p l e y L. S., "Proc. Nat. Acad. Sci.», 1953, v. 39, p. 1095-1100; [2] G i l l e t te D., B KH.: Contributions to the theory of games, v. 3, Princeton (N. J.), 1957, p. 179-87.

B. К. Доманский.



СТОХАСТИЧЕСКАЯ МАТРИЦА - квадратная (возможно, бесконечная) матрица $P=\left\|p_{i j}\right\| \mathrm{c}$ неотрицательными элементами такими, что



\[

\sum_{j} p_{i j}=1 \text { при любом } i \text {. }

\]





\end{document}